\section{Discussion}
\label{sec:discussion}

\subsection{Scope and Limitations}

The setup described in this article is explicitly scoped to the local development environment. Aspire's AppHost orchestrator is a development-time tool: it launches processes and containers on the developer's machine, injects configuration, and provides a telemetry dashboard. It is not a production deployment mechanism. Shipping the AppHost to a server and running it there is neither the intended use case nor a supported scenario.

For production deployment, the canonical path is to use the Aspire manifest --- generated by \texttt{dotnet run --project HobbitGame.AppHost -- --publisher manifest --output-path ./manifest.json} --- as input to a deployment tool such as \texttt{azd} (Azure Developer CLI) or \texttt{aspir8} (Kubernetes manifests). These tools translate the Aspire resource graph into cloud-native deployment artifacts. The Aspire documentation covers this path in detail~\cite{aspiredocs2024}.

\subsection{EnsureCreatedAsync vs. MigrateAsync}

The HobbitGame POC uses \texttt{EnsureCreatedAsync} for schema creation. This is appropriate when the schema is stable and no migration history is maintained. It has two limitations: (i) it cannot apply incremental schema changes --- if the schema is modified after the database exists, \texttt{EnsureCreatedAsync} is a no-op and the change is not applied; (ii) it is incompatible with EF Core migrations, because it creates the schema directly from the model and does not create the \texttt{\_\_EFMigrationsHistory} table. Teams intending to evolve the schema should add a migration (\texttt{dotnet ef migrations add Init}) and switch to \texttt{MigrateAsync} at startup.

\subsection{ContainerLifetime.Persistent in Team Environments}

\texttt{ContainerLifetime.Persistent} is designed for the single-developer local workflow. In a team environment where multiple developers share a machine (uncommon) or where CI runs the AppHost (not a standard CI pattern), the persistent lifetime may produce unexpected behaviour: one developer's container state may conflict with another's, or a CI agent may accumulate stale containers across runs. For these environments, the default \texttt{ContainerLifetime.Session} --- which tears down containers on AppHost exit --- is safer.

\subsection{Generalisation to Other Databases and Services}

The pattern described for PostgreSQL generalises directly to other databases and services supported by Aspire's hosting packages. \texttt{Aspire.Hosting.Redis} provides \texttt{AddRedis}; \texttt{Aspire.Hosting.SqlServer} provides \texttt{AddSqlServer}; \texttt{Aspire.Hosting.RabbitMQ} provides \texttt{AddRabbitMQ}. All follow the same pattern: declare the resource in the AppHost, add a reference to the dependent project, call the corresponding integration package's registration method in the dependent project's \texttt{Program.cs}. The secrets, health-check ordering, and telemetry mechanisms work identically across resource types.

\subsection{Relation to Docker Compose}

Teams already using Docker Compose for local development may ask whether Aspire replaces it. The honest answer is: for the local development scenario it largely does, and with a better developer experience (typed C\# builder calls, automatic connection-string injection, integrated dashboard, no YAML). However, Compose files remain useful for (i) sharing a reproducible environment with non-.NET tooling, (ii) running the stack in CI environments where the AppHost pattern does not apply, and (iii) deploying to Docker Swarm. The two approaches can coexist: a Compose file can be maintained for CI and non-development environments while the Aspire AppHost is used for local development.
